(1) Expand bilinearly. The $\pi_1\wedge\pi_2$ term is $N_0$. The cross terms give $\varphi(A)$, manifestly in $\Som\wedge\Spi$. The top term is $A\pi_1\wedge A\pi_2$, a bivector of the $2$-plane $\Som$, hence a multiple of $e_1\wedge e_2$; writing $A\pi_i = \sum_j A_{ji}e_j$, the multiple is $A_{11}A_{22} - A_{21}A_{12} = \det A$.
(2) $P(A)$ is decomposable by construction, hence null (Theorem \claimref{correspondence.spinor-identification}(2)). Expanding its norm with Proposition \claimref{correspondence.skew-lines-splittings-hyperbolic-planes-one} (the mixed slot is orthogonal to the hyperbolic plane):
\[
0 \;=\; Q(P(A)) \;=\; Q(\varphi(A)) + 2\det A\,\ip{N_0}{N_\infty}
 \;=\; Q(\varphi(A)) + 2\det A ,
\]
which is the displayed norm identity, and equivalently $\det A = -\tfrac12Q(m(A))$; substituting back into (1) yields $P(A) = X(m(A))$ exactly, in the chart normalisation of Section \ref{sec:setting}. Injectivity of $\varphi$: if $\varphi(A) = 0$ then $P(A) = N_0 + (\det A)N_\infty$ is decomposable, but $(N_0 + cN_\infty)\wedge(N_0 + cN_\infty) = 2c\,N_0\wedge N_\infty$, so $\det A = 0$ and $P(A) = N_0$, forcing $L_A = L_0$ and $A = 0$. Since $\dim \operatorname{Hom}(\Spi,\Som) = 4 = \dim(\Som\wedge\Spi)$, $\varphi$ is an isomorphism.
(3) A $2$-plane $U$ satisfies $U \cap \Som = 0$ if and only if the projection $U \to \Spi$ along $\Som$ is injective, hence bijective, if and only if $U$ is the graph of a unique $A: \Spi \to \Som$; and $U \cap \Som = 0$ is exactly skewness of $\mathbb{P}(U)$ and $L_\infty$. For the quadric statement: for decomposable $P = u_1\wedge u_2$, $\ip{P}{N_\infty} = \varepsilon(u_1\wedge u_2\wedge e_1\wedge e_2)$ vanishes exactly when the four vectors are dependent, i.e.\ when the planes meet; and $\ip{P(A)}{N_\infty} = \varepsilon(N_0\wedge N_\infty) = 1$. Rescaling any quadric point with $\ip{P}{\nInf} \neq 0$ to pairing $1$ places it in the chart, uniquely.
(4) (i) and (ii) are Proposition \claimref{correspondence.skew-lines-splittings-hyperbolic-planes-one} applied to the pair $(\nInf, \nO)$; (iii) is Proposition \claimref{correspondence.reality-conjugation-twistor-form-factor}(5). (iv): by equivariance of $\Psi$, a spin transformation stabilises the line $L_\infty$ (the plane $\Som$) if and only if it stabilises the ray of the bivector $N_\infty = \Psi(\nInf)$, i.e.\ the ray of $\nInf$ in $V_\C$. That stabiliser is computed in the flat model of \cite{ButlerT6}: the parabolic with Levi factor $\mathfrak{so}(1,3)\oplus\R$ and nilradical the translations, imported here as Fact \claimref{correspondence.imported-stabilisers-flat-model} below. The Levi factor preserves both $\Som$ and $\Spi$ (it stabilises both rays), and the assignment of the two $\mathfrak{sl}(2,\C)$ factors of $\mathfrak{so}(1,3)_\C$ to the two planes is Lemma \claimref{correspondence.chiral-planes-summand-spaces}(1). The verification script checks all eleven parabolic basis elements preserve $\Som$, that the complementary four do not, and the chiral assignment itself.
